\section{Main Results}
\label{sec:theorems}

Throughout this section, $\sigma$ denotes the M\"obius parameter
$\sigma = \log|\lam| = \mathrm{Re}(s)$ unless otherwise qualified.  The
jet half-width is denoted $\ell_J$ (not $\sigma_J$) to avoid confusion.

\subsection{Wavenumber selection (Saturn)}

\begin{proposition}[Stationary wavenumber]\label{prop:wavenumber}
For the barotropic Rossby dispersion relation $c = U - \beta/k^2$ with
$k = n/R_{\mathrm{hex}}$, the stationary condition $c = 0$ gives
\begin{equation}\label{eq:nstar}
  n^* = R_{\mathrm{hex}}\sqrt{\beta / U_{\max}}.
\end{equation}
For Saturn ($R_{\mathrm{hex}} = 5.1 \times 10^7$~m,
$\beta = 1.46 \times 10^{-12}$~m$^{-1}$\,s$^{-1}$,
$U_{\max} = 120$~m/s): $n^* = 5.62$, whose nearest integer is $n = 6$.
\end{proposition}

\begin{proof}
The Rossby dispersion relation for a barotropic wave with zonal
wavenumber $k$ and zero meridional wavenumber is $c = U - \beta/k^2$.
Setting $c = 0$ (stationary wave): $U = \beta/k^2$, giving
$k = \sqrt{\beta/U}$.  With $k = n/R_{\mathrm{hex}}$:
\[
  n = R_{\mathrm{hex}}\,k = R_{\mathrm{hex}}\sqrt{\beta/U}.
\]
Substituting Saturn's parameters:
\[
  n^* = 5.1 \times 10^7 \times \sqrt{\frac{1.46 \times 10^{-12}}{120}}
  = 5.1 \times 10^7 \times 1.103 \times 10^{-7} = 5.62.
\]
The nearest integer is $n = 6$.  The next candidates ($n = 5$: $U^* = 152$~m/s;
$n = 7$: $U^* = 77$~m/s) are substantially farther from the observed
$U_{\mathrm{obs}} = 120$~m/s.
\end{proof}

\subsection{Stationarity equivalence}

\begin{theorem}[Stationarity equivalence]\label{thm:stationarity}
Within the WKB approximation near the jet center, if the Rayleigh--Kuo
potential satisfies $V(\rhostar) < 0$ (a condition verified for Saturn's
jet parameters), then the stationary Rossby wave eigenfunction has
$\sigma = 0$, i.e., $|\lam| = 1$.
\end{theorem}

\begin{proof}
Near the jet center, the Rayleigh--Kuo potential
$V(\rho) = n^2 - (\beta - U'')/U$ is approximately constant over a
region of width $\sim \ell_J$.  Within this region, the WKB approximation
gives $\psihat_n(\rho) \approx C\,e^{s\rho}$ with $s^2 = V(\rhostar)$
and $s$ constant to leading order.

Since $V \in \mathbb{R}$, we have $s^2 \in \mathbb{R}$.  Writing
$s = \sigma + i\alpha_\rho$ (where $\alpha_\rho$ is the WKB radial
oscillation rate):
\[
  \mathrm{Im}(s^2) = 2\sigma\alpha_\rho = 0.
\]
If $V(\rhostar) < 0$, the solution is oscillatory: $s = \pm i\alpha_\rho$
with $\alpha_\rho = \sqrt{-V} > 0$.  Since $\alpha_\rho \neq 0$, we
conclude $\sigma = 0$.

\emph{Verification of $V < 0$:}  For Saturn's Gaussian jet at the jet
center ($\rho = \rhostar$), $(\beta - U''(0))/U(0) \approx
\ell_J^{-2} + \beta/U_{\max} \approx 400 + 10^{-16} \approx 400$
(dominated by the jet curvature), while $n^2 = 36$.  Therefore
$V(0) = 36 - 400 = -364 < 0$.  The condition holds by a wide margin.

With $\sigma = 0$, the real part gives the stationary jet speed:
$U^* = (\beta - U''(\rhostar)) / [(\alpha_\rho^2 + n^2)\,e^{-2\rhostar}]$.
For Saturn: $U^* \approx 105$~m/s, consistent with
$U_{\mathrm{obs}} = 120.1$~m/s \citep{Antunano2015}
(12\% discrepancy from the barotropic approximation).
\end{proof}

\begin{remark}[The identification $\alpha = 2\pi/n$]\label{rem:alpha}
The M\"obius azimuthal parameter $\alpha$ in $\lam = e^{\sigma + i\alpha}$
is \emph{identified} with $2\pi/n$ from the $\mathbb{Z}_N$ symmetry of
the polygon: an $N$-gon with vertices at $\theta_k = 2\pi k/N$ has a
discrete rotational symmetry of angle $2\pi/N$, which is by definition
$\arg(\lam)$ for the rotation $z \mapsto e^{i\,2\pi/N} z$.

This is not a derivation from the M\"obius framework---it is a restatement
of the definition of the multiplier for a linear rotation.  The M\"obius
framework contributes the \emph{continuous} family parameterized by
$\sigma$ (the spiral-to-circle transition); the \emph{discrete}
wavenumber $n$ comes from the QGPV (Proposition~\ref{prop:wavenumber}).

The WKB radial oscillation rate $\alpha_\rho$ satisfying
$\alpha_\rho^2 = |V(\rhostar)|$ is a separate quantity.  This paper
uses $\alpha$ for the M\"obius azimuthal parameter and $\alpha_\rho$ for
the WKB rate.
\end{remark}

\subsubsection{Branch disconnection within WKB}

\begin{proposition}[Branch disconnection]\label{prop:protection}
Within the WKB approximation, the condition $\sigma = 0$ is preserved
under perturbations of the jet speed $\delta U$.  The elliptic
($\sigma = 0$) and hyperbolic ($\alpha = 0$) branches of the eigenvalue
equation are disconnected.
\end{proposition}

\begin{proof}
The potential $V(\rhostar; \delta U)$ is real for all real $\delta U$.
The eigenvalue equation $s^2 = V$ requires $s^2 \in \mathbb{R}$.
Writing $s = i\alpha + \delta s$ with $\delta s = p + iq$:
\[
  \mathrm{Im}(s^2) = 2p(\alpha + q) = 0.
\]
This gives two branches:
\begin{enumerate}
  \item \textbf{Elliptic} ($p = 0$): $\sigma = 0$ for all $\delta U$.
  \item \textbf{Hyperbolic} ($q = -\alpha$): $\sigma = \sqrt{V + \alpha^2}$,
    purely radial flow.
\end{enumerate}
These branches meet only at $V = 0$ ($\delta U/U^* = \alpha^2/n^2
\approx 0.03$ for Saturn) and are not continuously connected for
$V \neq 0$.
\end{proof}

\begin{remark}[Scope of the protection]\label{rem:protection-scope}
Proposition~\ref{prop:protection} establishes branch disconnection
\emph{within the WKB approximation}, which is valid to leading order in
$\ell_J$.  Higher-order corrections to the eigenvalue equation could in
principle connect the branches; we have not excluded this.  The physical
persistence of Saturn's hexagon across four decades is consistent with
the branch disconnection but is not a consequence of
Proposition~\ref{prop:protection} alone.
\end{remark}

\subsection{The M\"obius parameter from the boundary value problem}

\begin{remark}[Multiple $\sigma$-like quantities]\label{rem:sigmas}
The M\"obius parameter $\sigma = \log|\lam|$ should be distinguished from:
$\sigma_{\mathrm{mode}} = \mathrm{Re}(s)$ in the QGPV eigenfunction
(which is zero within WKB); $\ell_J$, the jet half-width in log-polar
units; and $\sigma_{\mathrm{geom}}$, the observational polygon distortion
measure defined in \cref{sec:sigma-geom}.  Only $\sigma$ (the M\"obius
parameter) and $\sigma_{\mathrm{geom}}$ (the observable) appear in the
cross-planetary comparison.  They measure different things:
$\sigma$ characterizes the mean flow geometry (spiral winding);
$\sigma_{\mathrm{geom}}$ measures deviation from a perfect $N$-gon.
\end{remark}

The M\"obius parameter $\sigma$ is determined from the three-region
boundary value problem through a two-step chain:
\begin{enumerate}
  \item The BVP matching (below, Theorem~\ref{thm:matching}) determines
    the hexagon amplitude $\varepsilon$.
  \item The hexagonal boundary geometry gives the velocity ratio
    $|u_\rho/u_\theta|_{\mathrm{rms}} = n\varepsilon/\sqrt{2}$, and
    $\sigma = \alpha \cdot |u_\rho/u_\theta|$.
\end{enumerate}

\subsection{Thomson vortex ring stability}

\begin{proposition}[Critical central vortex strength]\label{prop:critical-ratio}
For the $N = 6$ Thomson ring on the unit circle with central vortex
$\kappa_0$, the characteristic polynomial of the $12 \times 12$
Hamiltonian stability matrix (constructed symbolically in Mathematica) is
\begin{equation}\label{eq:char-poly}
  p(\lam) \propto \lam^2 \cdot
  \underbrace{\bigl(9 + 36\kappa_0/\kappa + 16\pi^2\lam^2\bigr)}_{\text{Factor A}} \cdot
  \underbrace{Q(\lam^2, \kappa_0/\kappa)^2}_{\text{Factor B (double)}},
\end{equation}
where $Q$ is a quartic in $\lam^2$ with perfect-square discriminant
$(9 - 12\kappa_0/\kappa)^2$.  The first instability arises from Factor~A:
\begin{equation}\label{eq:critical-ratio}
  \boxed{\left(\frac{\kappa_0}{\kappa}\right)_{\!\mathrm{crit}} = -\frac{1}{4}.}
\end{equation}
\end{proposition}

\begin{proof}
The Jacobian of the Hamiltonian point vortex equations in the co-rotating
frame, including ring--ring and ring--center interactions, is constructed
symbolically.  Its characteristic polynomial, computed by Mathematica's
\texttt{Det} function and verified numerically at multiple $\kappa_0$
values, factors as \eqref{eq:char-poly}.  The stability threshold from
Factor~A is $9 + 36\kappa_0/\kappa = 0$, giving
$\kappa_0/\kappa = -1/4$.
\end{proof}

Remarkably, $\kappa_{\mathrm{crit}} = -1/4$ is specific to $N = 6$: for
$N = 4$ the critical ratio is $-1/2$ (proven algebraically via
characteristic polynomial), and numerical computation gives
$\kappa_{\mathrm{crit}} = -1/2$ for $N = 3$ and $N = 5$ as well
(to within $10^{-3}$; algebraic proof for these cases is obstructed by
irrational vertex coordinates).  The hexagonal ring thus has enhanced
sensitivity to anticyclonic central perturbations compared to smaller
rings.

\paragraph{Spectral non-equivalence.}
The Thomson stability matrix and the Rayleigh--Kuo operator share
$\mathbb{Z}_6$ eigenvectors (discrete Fourier modes) but not eigenvalues.
Galerkin projection gives eigenvalue ratios varying from $2.74$ (mode~1)
to $0.14$ (mode~3).  The operators are \emph{not} spectrally equivalent;
they converge at $|\lam| = 1$ independently.

\subsection{Matched asymptotic amplitude (Saturn)}

\begin{theorem}[Hexagon amplitude]\label{thm:matching}
The matched asymptotic expansion connecting the inner loxodromic flow
to the outer Rossby wave through the PV step at $\rhostar$ predicts the
hexagon amplitude
\begin{equation}\label{eq:amplitude}
  \varepsilon = C \cdot \frac{\delta U}{U^*} \cdot
  \exp\!\bigl(-n\,|\rhostar|\bigr),
\end{equation}
where $\delta U = U_{\mathrm{obs}} - U^*$ and $C$ is determined by the
Fourier projection of the PV gradient onto the eigenfunction envelope:
\begin{equation}\label{eq:overlap-integral}
  C = \frac{1}{U^* n}\int_0^\infty
  \bigl(\beta - U''(\rho)\bigr)\,e^{-n\rho}\,\dd\rho.
\end{equation}
For the Gaussian jet $U(\rho) = U_{\max}\exp(-\rho^2/2\ell_J^2)$:
\begin{equation}\label{eq:matching-constant}
  C = \frac{U_{\max}}{U^*} - \frac{n\,\ell_J}{U^*}
      \sqrt{\frac{\pi}{2}}\,e^{n^2\ell_J^2/2}\,
      \mathrm{Erfc}\!\left(\frac{n\,\ell_J}{\sqrt{2}}\right)
      + \frac{\beta}{n\,U^*}.
\end{equation}
For Saturn ($n = 6$, $\ell_J = 0.05$, $U_{\max} = 120$~m/s,
$U^* = 105$~m/s): $C = 0.797$, giving $\varepsilon = 0.112$.
\end{theorem}

This matches the Cassini-observed meander amplitude
$\varepsilon_{\mathrm{obs}} = 0.112$, under a Gaussian jet approximation
(9.8\% RMS fit to the actual wind profile; see \cref{sec:discussion} for
the propagated uncertainty).

\subsection{Energy--enstrophy sign structure}
\label{sec:enstrophy}

\begin{proposition}[Energy maximum at the polygon]\label{prop:energy-max}
For $N \ge 3$ point vortices on a ring, consider the spiral deformation
\begin{equation}\label{eq:spiral-deformation}
  z_k(\sigma) = R\,e^{\sigma k/N}\,e^{2\pi ik/N},
  \qquad k = 0, \ldots, N-1,
\end{equation}
which at $\sigma = 0$ gives the regular $N$-gon and for $\sigma > 0$
stretches vertex $k$ outward by the factor $e^{\sigma k/N}$, spiraling
the configuration.  The Thomson interaction energy
\begin{equation}\label{eq:H-sigma}
  H(\sigma) = -\sum_{j<k} \ln\bigl|z_j(\sigma) - z_k(\sigma)\bigr|
\end{equation}
satisfies $H''(0) < 0$.
\end{proposition}

\begin{remark}[Nature of the deformation]
The deformation \eqref{eq:spiral-deformation} is \emph{not} the uniform
M\"obius action $z_k \mapsto \lam z_k$ (which scales and rotates the
entire polygon rigidly, leaving $H$ unchanged).  It is a
\emph{symmetry-breaking} deformation that assigns each vertex a different
radial stretch, converting the $N$-gon into a logarithmic spiral.  This
direction in configuration space is physically relevant: it is the
deformation produced by a loxodromic flow $w(z) = Az^s$ acting on fluid
parcels at different azimuthal positions over time.  However, we do not
claim this is the \emph{unique} relevant deformation direction; $H''(0)$
could have different signs along other directions in the $N$-dimensional
configuration space.  The result is that the polygon is an energy maximum
\emph{along the spiral deformation direction}, not necessarily globally.
\end{remark}

\begin{proof}
At $\sigma = 0$, the pairwise squared distance is
$|z_j - z_k|^2 = 4R^2\sin^2(\pi m/N)$ where $m = k - j$.
Define $F(u, m) = \ln[1 - 2\cos(2\pi m/N)\,e^{um} + e^{2um}]$ where
$u = \sigma/N$, so that $\ln|z_j - z_k|^2 = 2uj + \ln R^2 + F(u,m)$.
Differentiating $F$ twice:
\[
  F''(0, m) = \frac{m^2}{1 - \cos(2\pi m/N)}
  = \frac{m^2}{2\sin^2(\pi m/N)}.
\]
Since $\dd^2(2uj)/\dd u^2 = 0$, the only contribution to $H''$ is from $F$.
Each pair $(j,k)$ with separation $m$ contributes $F''(0,m)$, and there
are $(N - m)$ such pairs for $m = 1, \ldots, N-1$.  Converting
$\dd/\dd\sigma = (1/N)\,\dd/\dd u$:
\begin{equation}\label{eq:Hpp-formula}
  H''(0) = -\frac{1}{4N^2}\sum_{m=1}^{N-1}
  \frac{(N-m)\,m^2}{\sin^2(\pi m/N)}.
\end{equation}
Every factor in the sum is strictly positive: $(N-m) > 0$, $m^2 > 0$,
$\sin^2(\pi m/N) > 0$ for $m = 1, \ldots, N-1$.  Therefore the sum is
positive and $H''(0) < 0$.
\end{proof}

The regularized enstrophy $Z(\sigma) = \sum_{j<k}|z_j - z_k|^{-2}$
satisfies $Z''(0) > 0$ (verified computationally for $N = 3, \ldots, 9$).
The ratio $\mu = -Z''/H'' > 0$ for all $N$, corresponding to negative
temperature in the Onsager sense.

\begin{remark}[Scope of the energy maximum]
The curvatures $H''(0)$ and $Z''(0)$ are computed along the specific
one-parameter loxodromic deformation defined above.  In the full
$N$-dimensional configuration space, perturbations exist that decrease
$Z$ (random tests show 62--78\% of perturbations increase $Z$, not
100\%).  The energy maximum is \emph{local} and
\emph{direction-dependent}, not global.

While $H''(0) < 0$ is proven for all $N \ge 3$
(Proposition~\ref{prop:energy-max}), the enstrophy result $Z''(0) > 0$
remains computational.
\end{remark}
